% TOP_PROBLEM: {"problemId": "problem.vinogradov-conjecture-on-the-least-quadratic-non-residue", "canonicalTitle": "Vinogradov Conjecture on the Least Quadratic Non-Residue", "releaseRank": 356, "primaryDomain": "number_theory_arithmetic_geometry", "primaryDomainLabel": "Number theory & arithmetic geometry", "displayStatus": "open", "releaseStatus": "open"}
% !TeX program = lualatex
% ATTEMPT_STATUS: unresolved
\documentclass[11pt]{article}
\usepackage[margin=25mm]{geometry}
\usepackage{fontspec}
\setmainfont{Latin Modern Roman}
\usepackage{amsmath,amssymb,mathtools}
\usepackage{unicode-math}
\setmathfont{Latin Modern Math}
\usepackage{xurl,hyperref}
\hypersetup{colorlinks=true,urlcolor=blue}
\setcounter{secnumdepth}{0}
\setlength{\parindent}{0pt}
\setlength{\parskip}{5pt}
\setlength{\emergencystretch}{3em}
\begin{document}
\section*{\texorpdfstring{Vinogradov Conjecture on the Least Quadratic Non-Residue}{Vinogradov Conjecture on the Least Quadratic Non-Residue}}\label{356}
\subsection{Definitions and mathematical statement}
For an odd prime $p$, a quadratic nonresidue is an integer $a$ not divisible by $p$ for which $x^2\equiv a\pmod p$ has no solution. Let $n_p$ be the least positive such integer. Vinogradov's conjecture in the selected form asserts
\[
 \forall{\varepsilon}>0\ \exists C_{\varepsilon}>0\ \forall p\text{ odd prime},
 \qquad n_p\le C_{\varepsilon} p^{\varepsilon}.
\]
The constant depends only on ${\varepsilon}$, not on $p$. Equivalently, $n_p=p^{o(1)}$ as primes tend to infinity. The requirement is unconditional, without assuming a Riemann hypothesis. No computable dependence of $C_{\varepsilon}$ on ${\varepsilon}$ or uniform numerical bound is demanded by this notation.
\par\smallskip\textit{Formulation and concept references:} \href{https://arxiv.org/pdf/2404.08380}{[S1]}.

\subsection{Short English statement}
Is the least quadratic nonresidue modulo a prime smaller, up to a constant, than every fixed positive power of that prime?
\subsection{Research attempt}
\textbf{Status: UNRESOLVED.} The unconditional bound $n_p=p^{o(1)}$ is not proved. The attempt gives elementary structural facts, derives a square-root character-sum bound, and explains in detail the smooth-number step which improves a Burgess input to the classical exponent $1/(4\sqrt e)$. The latter remains a positive exponent, and the much stronger logarithmic estimate in the cited source is conditional on GRH.

\paragraph{Source audit and exact target.}
Let $\chi$ be the quadratic character modulo the odd prime $p$, extended by $\chi(0)=0$, and let $n_p$ be the least positive integer with $\chi(n_p)=-1$. The goal quantifies over every $\varepsilon>0$, with an allowed constant $C_\varepsilon$ depending on $\varepsilon$ but not $p$. A single fixed positive exponent, however small, is insufficient. In [S1], the least-nonresidue estimate of logarithmic-square size is explicitly obtained under the Generalized Riemann Hypothesis. Its Fourier optimization does not remove that hypothesis. I checked this qualification in the primary manuscript on 27 September 2026.

\paragraph{1. Elementary facts and a special square-root bound.}
Exactly half the nonzero residues modulo $p$ are squares, since the map $a\mapsto a^2$ on $\mathbb F_p^\times$ is two-to-one. Thus $n_p$ exists and
\[
 2\le n_p\le\frac{p+1}{2}.
\]
The upper bound follows because the integers $1,\ldots,(p+1)/2$ are distinct nonzero residues and cannot all belong to a set of size $(p-1)/2$.

The integer $n_p$ is prime. Otherwise write $n_p=ab$ with $1<a,b<n_p$. Minimality gives $\chi(a)=\chi(b)=1$, contradicting $\chi(n_p)=\chi(a)\chi(b)=-1$. More generally, every positive integer whose prime factors are below $n_p$ and which is not divisible by $p$ is a quadratic residue.

For primes $p\equiv1\pmod4$ there is a short additional argument. Suppose $n=n_p>\sqrt p$, and divide $p=qn+r$ with $0<r<n$. The remainder is nonzero because $n$ is a prime smaller than $p$. Also $1\le q<n$. Thus $\chi(q)=\chi(r)=1$, while $r\equiv-qn\pmod p$ gives
\[
 1=\chi(r)=\chi(-1)\chi(q)\chi(n)=-1,
\]
since $\chi(-1)=1$ in this congruence class. Hence $n_p<\sqrt p$ when $p\equiv1\pmod4$. For $p\equiv3\pmod4$ the same calculation is consistent, rather than contradictory; dropping its sign would be an error.

The supplementary quadratic-reciprocity law gives $\chi(2)=-1$ for $p\equiv3,5\pmod8$, and then $n_p=2$. These elementary classes leave infinitely many primes untreated by a uniform small bound, and do not approach an arbitrary $p^\varepsilon$ estimate.

\paragraph{2. Derive a character-sum bound by finite Fourier analysis.}
Put $e_p(t)=\exp(2\pi it/p)$ and
\[
 \tau=\sum_{a\bmod p}\chi(a)e_p(a).
\]
For $t\not\equiv0\pmod p$, changing variables $a\mapsto t^{-1}a$ gives
\[
 \sum_{a\bmod p}\chi(a)e_p(ta)=\chi(t)\tau.
\]
The same Fourier coefficient is zero for $t=0$, because the nonzero residues split equally into squares and nonsquares. Parseval's identity for the finite Fourier transform therefore gives
\[
 (p-1)|\tau|^2
 =p\sum_{a\bmod p}|\chi(a)|^2=p(p-1),
 \qquad |\tau|=\sqrt p.
\]
Fourier inversion expresses the interval sum $S(H)=\sum_{n=1}^{H}\chi(n)$, for $1\le H<p$, as a sum of these coefficients against geometric progressions. For $1\le t\le p-1$,
\[
 \left|\sum_{n=1}^{H}e_p(tn)\right|
 \le\frac1{|\sin(\pi t/p)|}
 \le\frac{p}{2\min(t,p-t)}.
\]
The final inequality follows from $\sin(\pi x)\ge2x$ on $0\le x\le1/2$. Hence
\[
 |S(H)|\le\sqrt p\sum_{j=1}^{(p-1)/2}\frac1j
 \le\sqrt p\,(1+\log p).
\]
For $H=n_p-1$, every summand is one, so
\[
 n_p\le1+\sqrt p\,(1+\log p).
\]
This proof identifies exactly where the square-root scale enters: the nontrivial Fourier coefficients have magnitude $\sqrt p$, and the interval transform has a harmonic tail. It does not provide cancellation between those Fourier terms on much shorter intervals.

\paragraph{3. The Burgess input and the quarter-power barrier it crosses.}
Use the classical Burgess character-sum estimate as an external analytic theorem. In a form sufficient here, for a fixed positive integer $r$ and any fixed $\eta>0$,
\[
 |S(H)|\ll_{r,\eta}
 H^{1-1/r}p^{(r+1)/(4r^2)+\eta},\qquad 1\le H\le p.
\]
The harmless $p^\eta$ absorbs logarithmic factors; no proof of Burgess is claimed in this attempt. If $H$ is of order $p^\alpha$ with $\alpha>1/4$, divide by $H$. The power of $p$ becomes
\[
 -\frac\alpha r+\frac{r+1}{4r^2}+\eta.
\]
First choose $r$ so that $(r+1)/(4r)<\alpha$, and then choose $\eta$ sufficiently small. This power is negative, giving $S(H)=o(H)$ as $p\to\infty$.

If one used only $S(n_p-1)=n_p-1$, this would locate the least nonresidue at approximately the quarter-power scale. The standard improvement uses multiplicativity to create many additional known residues beyond $n_p$. The next paragraph makes that step explicit and tracks its limitation.

\paragraph{4. Smooth numbers and the factor $\sqrt e$.}
Assume for contradiction that every positive integer at most $y$ is a residue. Let $\Psi(x,y)$ count the positive integers at most $x$ whose prime factors are at most $y$. For $x<p$, all these integers have character one. Every other integer in the range has character at least minus one, so
\[
 S(\lfloor x\rfloor)\ge2\Psi(x,y)-\lfloor x\rfloor.
\]
Fix $1<u<2$ and set $x=y^u$, with $y\to\infty$. An integer at most $x$ can have at most one prime factor greater than $y$, counted with multiplicity, because two such factors would have product exceeding $y^2>x$. Thus the exact counting identity is
\[
 \Psi(x,y)=\lfloor x\rfloor-
      \sum_{y<q\le x\atop q\ \mathrm{prime}}\left\lfloor\frac{x}{q}\right\rfloor.
\]
The floor error is at most $\pi(x)=o(x)$. The standard reciprocal-prime asymptotic gives
\[
 \sum_{y<q\le y^u}\frac1q=\log u+o(1),
 \qquad \frac{\Psi(y^u,y)}{y^u}=1-\log u+o(1).
\]
The prime-number and reciprocal-prime estimates used here are classical external inputs; the combinatorial counting identity is exact and has been proved above. Consequently,
\[
 \frac{S(\lfloor y^u\rfloor)}{y^u}
 \ge1-2\log u+o(1).
\]
This lower bound stays positively separated from zero when $u<\sqrt e$.

Now let $\theta>1/(4\sqrt e)$. It suffices to consider $\theta<1/4$, since a bound with a smaller exponent implies one with a larger exponent after adjusting constants. Choose a fixed $u$ with
\[
 1<u<\sqrt e,\qquad \theta u>1/4.
\]
Put $y=\lfloor p^\theta\rfloor$ and $x=y^u$. Then $x<p$ for large $p$ and $x=p^{\theta u+o(1)}$. If $n_p>y$, the smooth-number argument gives $S(\lfloor x\rfloor)\ge c_u x$ for some $c_u>0$ and all sufficiently large $p$. Burgess gives $S(\lfloor x\rfloor)=o(x)$, a contradiction. Thus
\[
 n_p\ll_\theta p^\theta
 \quad\text{for every }\theta>\frac{1}{4\sqrt e}.
\]
The finitely many small primes are absorbed into the constant. This reconstructs the classical exponent from a named Burgess input; it is not a new result. The two strict inequalities on $u$ explain why it does not yield exponent zero: one needs enough smooth integers to force a positive character average and an interval long enough for the available cancellation estimate.

\paragraph{5. Least nonresidues cannot be bounded by one absolute integer.}
Fix an integer $H\ge2$ and let
\[
 Q=8\prod_{3\le q\le H\atop q\ \mathrm{prime}}q.
\]
By Dirichlet's theorem on primes in arithmetic progressions, there are infinitely many primes $p\equiv1\pmod Q$. For all such primes exceeding $H$, the supplementary law makes 2 a residue. For every odd prime $q\le H$, quadratic reciprocity and $p\equiv1\pmod4$ give
\[
 \left(\frac qp\right)=\left(\frac pq\right)
 =\left(\frac1q\right)=1.
\]
Multiplicativity then makes every integer $1\le n\le H$ a residue, so $n_p>H$. This proves that the least nonresidue is unbounded as $p$ ranges over primes, using Dirichlet's theorem as external input.

There is no conflict with $n_p=p^{o(1)}$. The argument chooses a possibly enormous prime after fixing $H$ and gives no lower bound of the form $n_p\ge p^\varepsilon$ for a fixed $\varepsilon>0$. Interchanging these quantifiers would turn a correct elementary construction into a false refutation of the target.

\paragraph{6. Exact diagnostics and the remaining analytic gap.}
The offline checks enumerate odd primes through $200000$, find their least nonresidues by Euler's criterion, and verify that each least nonresidue is prime. For the smaller primes through $2000$, direct enumeration of the nonzero square residues provides a second implementation. The special bound $n_p^2<p$ for $p\equiv1\pmod4$ is checked by integers. Character autocorrelation identities are checked for primes through 101, and smooth-number counts are compared with the exact large-prime-factor identity in ranges with $x<y^2$. These are finite arithmetic checks; they do not numerically establish a uniform exponent or GRH.

To complete the conjecture one needs a mechanism forcing a nonresidue below $C_\varepsilon p^\varepsilon$ for every positive $\varepsilon$. The Fourier estimate has a square-root scale. Burgess plus the explicit smooth-number argument reaches a smaller but fixed positive exponent. The source's logarithmic-square estimate depends on GRH. Removing these restrictions is exactly the missing work; no step here supplies the requisite unconditional short-interval cancellation or an alternative argument of comparable strength.

\subsection{Sources}
\begin{itemize}
\item[S1]E. Carneiro, M. B. Milinovich, E. Quesada-Herrera and A. P. Ramos, "Fourier optimization, the least quadratic non-residue, and the least prime in an arithmetic progression", arXiv:2404.08380 [math.NT], version dated 13 August 2025.
\url{https://arxiv.org/pdf/2404.08380}.
\textit{Formulation source.}
\end{itemize}
\end{document}
